\subsection{Régularité du maillage} \label{sec:regularite}

\noindent De la même manière que dans \cite{andreianov_discrete_2009}, on définit la \defini{taille} du maillage $\ensddfvplat$ par
\begin{equation} \label{eq:size_ensddfvplat}
\size(\ensddfvplat) \defeq \max_{E \in \ensprimalplat \cup \ensdualplat \cup \ensdiamantplat} \diam{E}.
\end{equation}
Pour mesurer l'applatissement des diamants, on note $\angle{\ensddfvplat}$ l'unique réel dans $\intervalleof{0}{\frac{\pi}{2}}$ tel que $\sin \angle{\ensddfvplat} \defeq \min_{\diamantplat{D} \in \ensdiamantplat} \abs{\sin \angle{\diamantplat{D}}}$. On introduit un nombre positif $\reg(\ensddfvplat)$ qui mesure la régularité d'un maillage donné et qui est utilisé pour démontrer la convergence des schémas :
\begin{multline} \label{eq:reg_ensddfvplat}
\reg(\ensddfvplat) \defeq \max\mathopen{}\Bigg\{
    \frac{1}{\sin \angle{\ensddfvplat}},
    \mathcal{N}, 
    \dual{\mathcal{N}},
    \max_{\diamantplat{D}_{\sigma, \dual{\sigma}} \in \ensdiamantplat}\mathopen{}\left[
    \frac{\mesure(\sigma)}{\mesure(\dual{\sigma})},
    \frac{\mesure(\dual{\sigma})}{\mesure(\sigma)}
    \right],
    \max_{\primalplat{K} \in \ensprimalplat} \frac{\diam{\primalplat{K}}}{\sqrt{\mesure(\primalplat{K})}},
    \max_{\mailledualeplate{K} \in \ensdualplat} \frac{\diam{\mailledualeplate{K}}}{\sqrt{\mesure(\mailledualeplate{K})}}, \\
    \max_{\primalplat{K} \in \ensprimalplat} \max_{\diamantplat{D} \in \ensdiamantplat_{\primalplat{K}}} \mathopen{}\left[ 
        \frac{\diam{\primalplat{K}}}{\diam{\diamantplat{D}}},
        \frac{\diam{\diamantplat{D}}}{\diam{\primalplat{K}}}
        \right],
    \max_{\mailledualeplate{K} \in \ensdualplat} \max_{\diamantplat{D} \in \ensdiamantplat_{\mailledualeplate{K}}} \mathopen{}\left[
        \frac{\diam{\mailledualeplate{K}}}{\diam{\diamantplat{D}}},
        \frac{\diam{\diamantplat{D}}}{\diam{\mailledualeplate{K}}}
        \right]
     \Bigg\},
\end{multline}
où $\mathcal{N}$ et $\dual{\mathcal{N}}$ sont le nombre maximum d'arêtes d'une maille primale et le nombre maximum d'arêtes incidentes d'un sommet du maillage primal. Ce nombre $\reg(\ensddfvplat)$ mesure essentiellement l'applatissement des diamants et le rapport entre le taille des mailles primales (resp. duales) et la taille des diamants et il doit être uniformément borné quand $\size(\ensddfvplat) \to 0$ pour avoir la convergence.

\begin{remarque}
Les quantités géométriques présentes dans la régularité sont celles des éléments plats, et qui sont donc calculables en pratique. 
\end{remarque}
Le \Cref{table:utilisation_termes_reg} donne, pour chaque terme de $\reg(\ensddfvplat)$, les résultats pour lesquels ils sont utilisés et les articles dans lesquels on les retrouve.
\begin{table}[H]
    \centering
    \begin{tblr}{
        % width=0.8\linewidth,
        % colspec={X[c]X[l]X[c]},
        colspec={ccc},
        hline{1,Z} = {1pt},
        row{2} = {abovesep+=5pt}
        }
    $\displaystyle\max_{\primalplat{K} \in \ensprimalplat} \max_{\diamantplat{D} \in \ensdiamantplat_{\primalplat{K}}} \frac{\diam{\primalplat{K}}}{\diam{\diamantplat{D}}}$, $\displaystyle\max_{\mailledualeplate{K} \in \ensdualplat} \max_{\diamantplat{D} \in \ensdiamantplat_{\mailledualeplate{K}}} \frac{\diam{\mailledualeplate{K}}}{\diam{\diamantplat{D}}}$ & \Cref{lemme:minoration_degre} & \cite{krell_schemas_2010} \\
    $\max\limits_{\diamantplat{D}_{\sigma, \dual{\sigma}} \in \ensdiamantplat}\mathopen{}\left[
        \displaystyle\frac{\mesure(\sigma)}{\mesure(\dual{\sigma})},
        \displaystyle\frac{\mesure(\dual{\sigma})}{\mesure(\sigma)}
    \right]$ & \Cref{prop:majoration_energie} & \cite{krell_schemas_2010} \\
    $\displaystyle\frac{1}{\sin \angle{\ensddfvplat}}$ & \Cref{prop:majoration_energie} & \cite{krell_schemas_2010} \\
    $\displaystyle\max_{\primalplat{K} \in \ensprimalplat} \max_{\diamantplat{D} \in \ensdiamantplat_{\primalplat{K}}} \frac{\diam{\diamantplat{D}}}{\diam{\primalplat{K}}}$, $\displaystyle\max_{\mailledualeplate{K} \in \ensdualplat} \max_{\diamantplat{D} \in \ensdiamantplat_{\mailledualeplate{K}}} \frac{\diam{\diamantplat{D}}}{\diam{\mailledualeplate{K}}}$ & \Cref{prop:mino_mesure_primal_K} & \cite{andreianov_discrete_2009} \\
    $\max\limits_{\primalplat{K} \in \ensprimalplat} \displaystyle\frac{\diam{\primalplat{K}}}{\sqrt{\mesure(\primalplat{K})}}$, $\max\limits_{\mailledualeplate{K} \in \ensdualplat} \displaystyle\frac{\diam{\mailledualeplate{K}}}{\sqrt{\mesure(\mailledualeplate{K})}}$ & \Cref{prop:mino_mesure_primal_K} & \cite{krell_schemas_2010} \\
    $\mathcal{N}$, $\dual{\mathcal{N}}$ & \Cref{prop:inegalites_strategie_un} & \cite{krell_schemas_2010} \\
    \end{tblr}
    \caption{Utilisation des termes de $\reg(\ensddfvplat)$.}
    \label{table:utilisation_termes_reg}
\end{table}
